\documentclass{elsarticle}
\usepackage[utf8]{inputenc}
\usepackage[T1]{fontenc}
\usepackage{geometry}
\usepackage{graphicx}
\usepackage{float}
\usepackage{hyperref}
\usepackage[english,indonesian]{babel}

\title{Optimasi Pemrosesan Template LaTeX Menggunakan Pendekatan Berbasis AI}
\author{Academic Support Agent}
\date{\today}

\begin{document}
\maketitle

\begin{abstract}

\end{abstract}

\chapter{PENDAHULUAN}


\chapter{TINJAUAN PUSTAKA}


% Original content (commented out):
% %% 
%% Copyright 2007-2020 Elsevier Ltd
%% 
%% Este arquivo é uma adaptação de arquivo do pacote 'Elsarticle', que está disponível respeitando-se os termos da licença 'LaTeX Project Public License', na versão 1.2 ou posterior. 
%%-------------------------------------------------------------


\documentclass[preprint,12pt]{elsarticle}

%% para incluir figuras, graphicx.sty foi adicionado ao arquivo
%% elsarticle.cls. Se você prefere comandos mais antigos
%% por favor, adicione o pacote \usepackage{epsfig}

%% O pacote amssymb contém cários símbolos matemáticos úteis
\usepackage{amssymb}

%% O pacote babel coloca acentuação e hifenização de acordo com o português (Brasil)
\usepackage[brazil]{babel}

%%usando padrão utf8 para codificação
\usepackage[utf8]{inputenc}

%% O pacote amsthm contém a definição de alguns ambientes matemáticos
\usepackage{amsthm}

%% O pacote lineno adiciona número às linhas do texto
\usepackage{lineno}

\journal{WPCCG}

\begin{document}

\begin{frontmatter}

\title{Título do seu Manuscrito}

%Resumo 
\begin{abstract}
%% Texto do resumo
\begin{resumop}
No meio do caminho tinha uma pedra tinha uma pedra no meio do caminho tinha uma pedra no meio do caminho tinha uma pedra. Nunca me esquecerei desse acontecimento na vida de minhas retinas tão fatigadas. Nunca me esquecerei que no meio do caminho tinha uma pedra tinha uma pedra no meio do caminho no meio do caminho tinha uma pedra (Carlos Drummond de Andrade)
\end{resumop}

\begin{resumoi}
Lorem ipsum dolor sit amet, consectetur adipiscing elit, sed do eiusmod tempor incididunt ut labore et dolore magna aliqua. Ut enim ad minim veniam, quis nostrud exercitation ullamco laboris nisi ut aliquip ex ea commodo consequat. Duis aute irure dolor in reprehenderit in voluptate velit esse cillum dolore eu fugiat nulla pariatur. Excepteur sint occaecat cupidatat non proident, sunt in culpa qui officia deserunt mollit anim id est laborum.
\end{resumoi}
\end{abstract}

\begin{keyword}
%% palavras-chaves devem ser colocadas aqui, no seguinte formato: palavra1 \sep palavra2 \sep palavra3
palavra um \sep palavra dois\sep palavra três
\end{keyword}

\end{frontmatter}

\linenumbers

%% testo princicpal
\section{Título de uma Seção Simples}
\label{sec:sample1}

%% Para citações use: 
%%       \citet{<label>} ==> Jones et al. [21]
%%       \citep{<label>} ==> [21]
%%

Lorem ipsum dolor sit amet, consectetur adipiscing \citep{Fabioetal2013} elit, sed do eiusmod tempor incididunt ut labore et dolore magna \citet{Blondeletal2008} aliqua. Ut enim ad minim veniam, quis nostrud exercitation ullamco laboris nisi ut aliquip ex ea commodo consequat. Duis aute irure dolor in reprehenderit in voluptate velit esse cillum dolore eu fugiat nulla pariatur. Excepteur sint occaecat cupidatat non proident, sunt in culpa qui officia deserunt mollit \citet{Blondeletal2008,FabricioLiang2013} anim id est laborum.

Lorem ipsum dolor sit amet, consectetur adipiscing elit, sed do eiusmod tempor incididunt ut labore et dolore magna aliqua. Ut enim ad minim veniam, quis nostrud exercitation ullamco laboris nisi ut aliquip ex ea commodo consequat. Duis aute irure dolor in reprehenderit in voluptate velit esse cillum dolore eu fugiat nulla pariatur. Excepteur sint occaecat cupidatat non proident, sunt in culpa qui officia deserunt mollit anim id est laborum see appendix~\ref{sec:another}.

\section{Another Section}
\label{sec:another}
Lorem ipsum dolor sit amet, consectetur adipiscing elit, sed do eiusmod tempor section \ref{sec:sample1} incididunt ut labore et dolore magna aliqua. Ut enim ad minim veniam, quis nostrud exercitation ullamco laboris nisi ut aliquip ex ea commodo consequat. Duis aute irure dolor in reprehenderit in voluptate velit esse cillum dolore eu fugiat nulla pariatur. Excepteur sint occaecat cupidatat non proident, sunt in culpa qui officia deserunt mollit anim id est laborum. Veja a  Figura~\ref{fig:net}.

\begin{figure}[h]
    \centering
    \includegraphics[width=0.25\textwidth]{net}
    \caption{Uma rede}
    \label{fig:net}
\end{figure}

%% Se você tem um arquivo bib e quer que o  bibtex gere os 
%% bibitems, por favor, use:
%%
\bibliographystyle{elsarticle-num-names} 
\bibliography{cas-refs} %no lugar de cas-refs coloque o nome do seu arquivo .bib

%% senão use o seguinte código para adicionar as entradas dos bibitems diretamente no arquivo TeX.

% \begin{thebibliography}{00}

% %% \bibitem[Autor(ano)]{label}
% %% Texto do item bibliográfico

% \bibitem[ ()]{}

% \end{thebibliography}
\end{document}

\endinput
%%
%% Fim o arquivo main.tex'.


\end{document}